% =============================================================================
% Section 1 -- Introduction
% =============================================================================
\section{Introduction}
\label{sec:introduction}

Organizations are increasingly utilizing conversational agents as their main
customer interfaces, managing tasks such as support tickets, returns, and
billing disputes at scale. The evaluation prior to release is now well-equipped
with standardized benchmarks~\cite{liang2022holistic}, red-teaming~\cite{ganguli2022redteaming,perez2022redteaming},
and reinforcement learning from human feedback (RLHF)
alignment~\cite{ouyang2022instructions}, creating a robust pre-deployment
lifecycle. However, there is a noticeable lack of a culture focused on
post-deployment remediation. A comprehensive study of deployed agents found
that many either fail or do not perform as expected, yet there is limited
understanding of how organizations react in these situations~\cite{pan2025measuring}.

Ad-hoc responses to degraded agents---including hallucinations about policies
cited, failure in a new product category, and changes in tone---are standard.
Typically, engineers will ``fix'' individual prompts, product managers will
create service requests (tickets), and occasionally the system will be rolled
back. There is no structure to how this occurs; there is no formalized way to
document and evaluate how well the fixes worked. Agents either recover through
informal means or they are quietly removed.

While this may seem like a minor operational issue, the lack of formal
remediation for degraded agents represents a substantial governance risk as AI
becomes more responsible for decisions with greater consequences, and
regulatory frameworks such as the European Union's AI Act require documentation
and human oversight of deployed AI systems~\cite{eu2024aiact}. Wachter et
al.~\cite{wachter2017counterfactual} have pointed out that traceability of the
decision-making process is required for accountability in automated decision-making
systems, which can be provided by formalized remediation processes. Furthermore, the
lack of formal remediation creates unnecessary regression loops; decreases
operator trust; and limits organizations' ability to determine when an agent
has failed to such a degree that it would be better to retire the agent.

We recommend using a performance improvement plan (PIP) as a model for
structuring remedial actions; we suggest using the PIP because it is an
established tool in organizational HR and can be used as a useful template for
remedying underperforming agents. A PIP includes a number of elements which are
common to all accountability protocols: a triggering condition, a means of
identifying the causes of failure, a limited time frame within which corrective
action can take place, regular check-ins to monitor progress toward goals and/or
compliance with requirements, and an explicit definition of when the subject has
sufficiently recovered to allow reintegration into normal operations or when it
will no longer be able to continue performing the required job functions.
Additionally, in the field of organizational behavior, Kirkpatrick's~\cite{kirkpatrick1994evaluating}
four-level model for evaluating training programs---reaction, learning,
behavior, results---is a useful complement to the above-mentioned PIP model,
providing a framework by which one can evaluate if the remedial actions taken
were successful in producing lasting positive changes in both individual and
aggregate system behaviors.

We make the following contributions:
\begin{itemize}
  \item We introduce the Agent Performance Improvement Plan (APIP) as a
        formalized post-deployment remediation lifecycle protocol, grounded in
        organizational management and training evaluation research.
  \item We develop a five-category taxonomy of agent failure modes that maps
        onto APIP cause-finding procedures, together with a published
        deterministic mapping from trace-level failure annotations
        (MAST~\cite{cemri2025mast}) and disaggregation signatures onto taxonomy
        categories and intervention tiers.
  \item We propose a three-tier intervention hierarchy---prompt-level,
        retrieval-level, and model-level---with escalation criteria, and connect
        the evaluation of these interventions to Kirkpatrick's multi-level
        evaluation framework.
  \item We formalize the APIP as a declarative schema to support tooling
        integration and audit trail generation.
  \item We evaluate the full protocol in three live multi-agent environments
        spanning customer service, software engineering, and economic
        operations, scoring detection lag, attribution accuracy (as a confusion
        matrix against injected ground truth, with abstention as a scored
        outcome), and oracle-relative time-to-resolution against a four-condition
        baseline family. The same contract/trigger/tier/exit machinery governs
        all three domains, which is itself evidence for the generality of the
        behavioral contract abstraction.
  \item We implement the mesh disruption problem---where a locally optimal new
        agent is hypothesised to degrade peer performance---together with a
        scoped-context onboarding protocol, and report a null result that
        localises the phenomenon's precondition: a bounded shared context.
\end{itemize}
